\documentclass[../../book-main.tex]{subfiles}
\graphicspath{{\subfix{../..}}}

\begin{document}

\chapter{Future Study of Intelligence}
\label{ch:future}

% \begin{quote}
% ``{\em If I were to choose a patron saint for cybernetics
% out of the history of science, I should have to choose
% Leibniz. The philosophy of Leibniz centers about two
% closely related concepts -- that of a universal symbolism
% and that of a calculus of reasoning.}''

% $~$\hfill -- Norbert Wiener, Cybernetics, 1961
  
% \end{quote}

\begin{quote}
``{\em The study is to proceed on the basis of the conjecture that every aspect of learning or any other feature of intelligence can in principle be so precisely described that a machine can be made to simulate it. An attempt will be made to find how to make machines use language, form abstractions and concepts, solve kinds of problem now reserved for humans, and improve themselves.}''

$~$\hfill -- Proposal for the Dartmouth AI program, 1956
 \end{quote}
\vspace{5mm}

%This chapter discusses active ongoing research topics on intelligence as well as some challenges and opportunities for future study.

Generally speaking, this manuscript is meant to systematically introduce mathematical principles and computational mechanisms behind how memory or knowledge can be developed from empirical observations. The capability to seek parsimony in a seemingly random world is a fundamental characteristic of any intelligence, natural or man-made. We believe that the principles and mechanisms presented in this book are rather unifying and universal and are applicable to both animals and machines. 

We hope that this book can help the readers fully clarify the mystery around modern practices of artificial deep neural networks by developing a rigorous understanding of their functions and roles in serving the objective of learning low-dimensional data distributions. With such a understanding, we should have become aware of both capabilities and limitations of existing AI models and systems: 
\begin{enumerate}
    \item Existing realizations of these functions are still rather primitive and certainly far from optimal in terms of optimization techniques or strategies. 
    \item Existing models and systems are short of being complete in terms of a memory system that is capable of self-learning and self-improving.
    \item Existing models and methods are for memorizing knowledge with empirical observations, but are not yet capable of developing human-level intelligence.
\end{enumerate} 

One of the goals of this book is for people to establish an objective understanding of current machine intelligence technologies and to realize what open challenges remain ahead for the study of intelligence. In the last chapter of the book, we provide some of our views and projections for the future. 

\section{Towards Autonomous Intelligence: Close the Loop?}
From the practice of machine intelligence in the past decade, it has become clear that, if there were sufficient data and computational resources, one can build a large enough model and train it end-to-end to learn the distribution of all the data. Such a pre-trained large model can even demonstrate superhuman capabilities in terms of memorizing existing knowledge about the world\footnote{For example,  those knowledge already encoded in massive texts.}. In a way, such large models play a similar role as DNAs for the lives to record and pass on knowledge about the world.   

Nevertheless, as we also know now, such open-ended models themselves do not have the mechanisms to verify whether the knowledge they have memorized from the data are correct, nor have the ability to revise or improve them in case of being incorrect or incomplete. Improvements in these models depend on human involvement: experimentation, evaluation, and selection.\footnote{We may call this process as ``artificial selection,'' as opposed to the natural selection for the evolution of the lives.} 

As we have studied earlier in this book (Chapter \ref{ch:closed-loop} in particular), closed-loop systems have the ability to align their internal representations with the (sensed) observations of the external world and continue to improve. An immediate step forward for the future is to develop and build truly closed-loop memory systems that are capable of learning and improving acquired knowledge autonomously and continuously, as shown in Figure \ref{fig:open-to-closed}.

To a large extent, we believe that transition from open-loop end-to-end trained models to closed-loop continuously learning systems:
\begin{equation}
   \mbox{\textbf{open-ended} models} \; \Longrightarrow \; 
   \mbox{\textbf{closed-loop} systems}
\end{equation}
is the key for machines to truly emulate how the (animal) brain creates and stores knowledge in an open world: 
\begin{quote}
\begin{center}
        {\em Open-ended models are for a closed world, however large; \\ closed-loop systems are for an open world, however small.}
\end{center}
\end{quote}

\begin{figure}[t]
    \centering    \includegraphics[width=0.9\linewidth]{figs_chap8/open-to-closed.png}
    \caption{From open-ended deep networks to closed-loop systems.}
    \label{fig:open-to-closed}
\end{figure}

The practice of artificial intelligence in the past few years has led many to believe that one needs to build one single large model, either an open-ended model or a closed-loop system, to memorize all knowledge. In the self-consistent setting, it is unreasonable to assume that the capacity of one closed-loop autoencoder would be sufficient. Even if this might be technologically possible, it is likely that such a solution is far from necessary and efficient. In addition, as we have alluded to in Chapter \ref{ch:autoencoding}, an important goal of learning a (compressive) autonencoder is to help {\em organize} the data distribution in terms of the codes and features learned. As single closed-loop autoencoder might not be the most suitable for this purpose. 

As we have learned from neuroscience, the cortex of our brain consists of tens of thousands of cortical columns. These cortical columns all have similar physical structures and functions. They are highly parallel and distributed, though sparsely interconnected. Hence, we believe that in order to develop a more scalable, interpretable, and modular memory system, we need to consider architectures that emulate that of the cortex. Figure \ref{fig:loops} shows a hypothesized architecture -- a massively distributed hierarchical system that consists of many largely parallel closed-loop autoencoding modules. These modules learn to  encode different sensing modalities or different aspects of information from each sensing modality. Higher-level autoencoders can be built on top of the lower-level ones to develop more sparse and lossy ``abstractions'' of the representations learned by the lower levels. 
\begin{figure}[t]
\centering
\includegraphics[width=0.99\linewidth]{figs_chap8/loops.png}
    \caption{Conjectured architecture of the brain cortex: The cortex is an auto-encoding system that consists of a hierarchy of closed-loop auto-encoders that extract information and maximize the information gain of the resulting representations at all levels.}
    \label{fig:loops}
\end{figure}


\section{Towards Intelligence of Nature: Beyond BP?}
As we have known from the practice of training deep networks, the only known scalable method to train such networks at scale is through back propagation (BP) \cite{Back-Prop}. Although BP has offered a way to correct errors via gradient signals propagated back through the whole model, it is nevertheless a method that differs significantly from how nature learns: BP is an option that nature cannot afford in terms of its high cost or simply cannot implement due to physical limitations. 

We cannot truly understand intelligence unless we also understand how it can be physically implemented efficiently. That is, one needs to address the computational  complexity of realizing mechanisms associated with achieving its objectives. Note that historically our understanding of (machine)  intelligence has gone through several phases in the past:
\begin{equation}
   \mbox{\textbf{incomputable}} \;
   \Longrightarrow \; \mbox{\textbf{computable}} \;
   \Longrightarrow \; \mbox{\textbf{tractable}} \; \Longrightarrow \; 
   \mbox{\textbf{scalable}}.
\end{equation}
To a large extent, the success of deep learning and back propagation is because they have enabled scalable implementation with modern computing platforms (such as GPUs). Nevertheless, such an implementation is still way too expensive for nature. 

We believe the future study of intelligence needs to enable the   following transition:
\begin{equation}
   \mbox{\textbf{scalable} solutions} \; \Longrightarrow \; 
   \mbox{\textbf{natural} solutions}.
\end{equation}
We need to discover new learning architectures and optimization mechanisms that enable learning under natural physical conditions and resource constraints, say, without conducting global optimization by BP or without access to data all at once. The new optimization strategies and mechanisms should work under conditions similar to those for intelligent beings in nature. They can be realized with significantly less resource requirements, amenable to continuous learning with localized and sparse update. These mechanisms may involve both backward and forward, continuous, local, distributed optimization. So far, the only case for which we know of such a solution is the simplest case of PCA, as seen in Chapter \ref{ch:autoencoding}. We believe that the hierarchical distributed closed-loop system architecture illustrated in Figure \ref{fig:loops} opens up many more possibilities than the single-model architecture. That makes such new learning mechanisms possible. 

\section{Towards Intelligence of Human: Turing Test?}
As we now know, humans have achieved two quantum leaps of intelligence in history. The first is the development of spoken and written languages about five to ten thousand years ago. That has enabled sharing and passing on knowledge, similar to the function of DNAs in nature. The second is the development of mathematics and logic about three thousand years ago. This has liberated us from summarizing knowledge from observations in empirical forms. They have become the language of science. We are now capable of developing and verifying new theories with abstract concepts such as numbers, space, time, and energy. 

As we have discussed in the Introduction (Chapter \ref{ch:intro}), the 1956 ``artificial intelligence'' (AI) program aimed to study high-level functions such as abstractions, logical inference, and problem solving that are believed to differentiate humans from animals: 
\begin{equation}
   \mbox{\textbf{low-level} (animal) intelligence} \; \Longrightarrow \; 
   \mbox{\textbf{high-level} (human) intelligence.}
\end{equation}
As we have clarified in this book, much of the technological advance in the past decades, carried out under the name ``AI'', is actually more closely related to the low-level intelligence shared by both animals and humans, such as forming memory. So far, there has been no evidence that suggests these mechanisms would suffice to achieve high-level intelligence that the original AI program truly cares about. 

In fact, we know little about how to rigorously verify whether a system truly has developed certain high-level intelligence, despite the fact that the Turing test has been proposed since 1950 \cite{Turing-1950}. For long, there was not deemed necessary since the capabilities of machines were far below that of a human. However, given recent technological advancements, many models and systems have claimed to reach and even surpass human intelligence. Therefore, it is high time to give a scientific and executable definition of the Turing test. That is, how do we systematically and objectively evaluate the level of intelligence for a given model or system? For example, how can we rigorously verify whether an intelligent system has truly grasped a certain abstraction concept, such as the notion of natural numbers, or merely memorizes a large number of its instances? Note that the state-of-the-art large language models still struggle with simple mathematical questions like ``{\em whether 3.11 is larger or smaller than 3.9?}''\footnote{Some large models have corrected their answers to this question via post engineering. However, we leave it to the reader as an exercise to test whether any of the state-of-the-art language models truly understand the notion of natural numbers and the associated arithmetic.} How to verify whether a system truly understands the rules of logic and knows how to apply them rigorously or has simply memorized a large number of instances of practicing logic? Furthermore, is such a system even capable of correcting its own knowledge or developing new knowledge such as physical laws, mathematical concepts, or causal relationships? In summary, it is high time that we develop rigorous evaluation methods that can tell a system/model's seemingly intelligent capability belongs to which of the following: 
\begin{enumerate}
    \item simply having memorized a large  amount of knowledge-carrying data and regenerating them;
    \item being able to autonomously and continuously develop new knowledge from new observations;
    \item actually having learned and understood abstract knowledge in texts and knowing how to apply correctly;
    \item being able to generate new logical, causal, or mathematical hypotheses or theories and verify them.
\end{enumerate}

As we have seen in this book, {\em compression} has played a fundamental role in learning. It is the governing principle and mechanism for identifying and organizing information encoded in an (empirical) data distribution. To a large extent, it can explain most of the practice of ``artificial intelligence'' in recent years. An outstanding question for future study is whether {\em compression alone} is sufficient to achieve all higher-level capabilities of intelligence listed above? Whether understanding, abstraction, causal inference, and logical reasoning are certain extreme or varied forms of compression? Is there a fundamental difference between compressing empirical observations and extracting certain high-level concepts and laws from empirical knowledge? We believe that uncovering and understanding the underlying mathematical principles and computational mechanisms of intelligence at different levels is the last frontier for science, mathematics, and computation together. 

\end{document}